%% SCRAP: papers/SCIENTIFIC_DEFENSE_CHECKLIST
%% SOURCE: docs/working/papers/SCIENTIFIC_DEFENSE_CHECKLIST.md
%% STATUS: CURRENT
%% FITS: ssrn/app-defense, vol3-research/ch-repro
%% EDITORIAL: lifted — prose rewritten to press voice; maturity matrix preserved

\section{Scientific Defense Checklist}
\label{sec:defense-checklist}

This section maps each plausible attack vector against the work, assesses
the strength of current defenses, identifies remaining gaps, and prescribes
specific remediation actions. The maturity matrix at the end
(\S\ref{sec:defense-matrix}) provides a one-table summary for submission
review.

\subsection{Attack Vector 1: Data Integrity}

\paragraph{Potential challenges.}
Fabricated data; cherry-picked runs; results too perfect; selective reporting.

\paragraph{Current defenses.}
Raw data for all 90 runs is committed to the repository under
\texttt{docs/archive/phase-1/Reference/physics\_experiment/PEER\_REVIEW\_SUBMISSION/03\_EXPERIMENTAL\_DATA/}.
Git history demonstrates iterative development rather than one-shot fabrication.
All 90 runs carry timestamps in \texttt{experiment\_summary.txt}.

\paragraph{Remaining gaps.}
SHA256 checksums for raw data files are not yet generated; GPG commit signing
for experimental data commits is pending; third-party verification by an
external researcher has not yet occurred; Zenodo DOI for permanent archival
is planned but not yet executed.

\paragraph{Remediation.}
\begin{lstlisting}[language=bash]
find docs/archive/phase-1/Reference/physics_experiment/ \
  -type f -exec sha256sum {} \; > EXPERIMENTAL_DATA_CHECKSUMS.txt
gpg --clearsign EXPERIMENTAL_DATA_CHECKSUMS.txt
# Then upload dataset to Zenodo for DOI
\end{lstlisting}

\subsection{Attack Vector 2: Statistical Validity}

\paragraph{Potential challenges.}
Misapplied ANOVA; insufficient sample size; $p$-hacking; multiple testing
without correction; incorrect null hypothesis.

\paragraph{Current defenses.}
ANOVA and Levene's test are documented in the formal claim table. Effect
sizes are reported alongside $p$-values (CV measurements, Cohen's $d$). The
DoE methodology was documented before experiments were run.

\paragraph{Remaining gaps.}
A formal power analysis justifying $N = 30$ is not yet written as a
standalone document. A reproducible R/Python statistical analysis script is
not yet committed. Normality tests (Shapiro-Wilk) and assumption validation
are not formally documented. Bootstrapped confidence intervals are not
reported.

\paragraph{Remediation (R script skeleton).}
\begin{lstlisting}[language=R]
library(tidyverse); library(car)
data <- read_csv("experimental_data.csv")
# Power analysis
power.t.test(n=30, delta=0.254, sd=0.05, sig.level=0.05)
# Normality
shapiro.test(data$performance)
# Variance homogeneity
leveneTest(performance ~ configuration, data=data)
# Effect size
cohen.d(early_runs, late_runs)
\end{lstlisting}

\subsection{Attack Vector 3: Experimental Design}

\paragraph{Potential challenges.}
Flawed methodology; uncontrolled confounds; invalid baseline comparison;
undocumented experimental conditions.

\paragraph{Current defenses.}
DoE methodology is documented. The baseline (C\_NONE: all loops off) is
defined. Workload and hardware were held constant.

\paragraph{Remaining gaps.}
Hardware specification (exact CPU model, RAM, kernel version) needs to be
formalized in a single canonical \texttt{EXPERIMENTAL\_PROTOCOL.md}. CPU
governor, thermal throttling, and ASLR settings are documented in the
reproducibility guide but not in a consolidated protocol document. Workload
entropy ($H = 3.2$\,bits, Zipf $\alpha = 1.1$) should be cited with
measurement provenance.

\paragraph{Protocol template.}
The experimental protocol should specify: CPU (Intel Xeon Gold 6154 @
3.00\,GHz, 18 cores); RAM (128\,GB DDR4-2666 ECC); OS (Ubuntu 22.04.3 LTS,
kernel 6.2.0-39-generic); GCC 11.4.0; CPU governor: performance; Turbo
Boost: disabled; ASLR: disabled; process affinity: core 0; 10-second
inter-run delay for thermal stabilization; timer source:
\texttt{CLOCK\_MONOTONIC\_RAW}; workload: Fibonacci(20), 10{,}000 iterations,
$\approx 2.1 \times 10^6$ word executions.

\subsection{Attack Vector 4: Claims versus Evidence}

\paragraph{Potential challenges.}
Claimed value differs from data; ``25.4\% improvement'' is misleading;
determinism claim is overstated; interpretation overreaches evidence.

\paragraph{Current defenses.}
Exact claims are enumerated in the formal claim table with data locations.
The qualifier ``algorithmic variance'' (not ``total variance'') is used
consistently.

\paragraph{Remaining gaps.}
Sensitivity analysis showing robustness of the 0\% claim across parameter
variations is provided separately (\S\ref{sec:sensitivity-analysis}) but
not yet wired into a single claim-to-evidence table. Limitations are not
formally listed in every document.

\subsection{Attack Vector 5: Reproducibility}

\paragraph{Potential challenges.}
Independent researcher cannot reproduce; build instructions incomplete;
missing dependencies; environment not fully specified.

\paragraph{Current defenses.}
Build instructions are in \texttt{README.md}. Source code is public.
780\raisebox{0.5ex}{+} tests validate correctness.

\paragraph{Remaining gaps.}
Docker container and one-command full experiment replay are planned but not
yet committed. CI/CD validation for automated reproduction regression
detection is planned (GitHub Actions workflow skeleton is documented in the
source).

\subsection{Attack Vector 6: Prior Art and Novelty}

\paragraph{Potential challenges.}
Work already done by JIT compilers; uncited important papers; incremental
not novel; contribution overstated.

\paragraph{Current defenses.}
Claims are specific: ``0\% algorithmic variance'' is not the same as
``adaptive JIT.'' The unique combination of determinism, adaptation, and
formal verification potential is identified.

\paragraph{Remaining gaps.}
A comprehensive literature review and a comparison table (this work versus
PyPy, HotSpot, LuaJIT, Chez Scheme, Ertl 1996) are missing. Key citations
to add: Bolz et al.\ (2009), Ertl (1996), Garlan et al.\ (2004).

\paragraph{Novelty statement.}
What is novel: deterministic adaptation (0\% algorithmic CV); Levene's-test-
driven variance-based window tuning; empirically validated steady-state
convergence; explicit thermodynamic metaphor mapping. What is not novel:
execution frequency tracking (standard profiler technique); exponential decay
(standard caching algorithm); hot-code acceleration (standard JIT technique).
The contribution is the \emph{combination}, not individual techniques.

\subsection{Attack Vector 7: Implementation Correctness}

\paragraph{Potential challenges.}
Code has bugs; implementation diverges from description; performance claims
from buggy code; tests miss critical paths.

\paragraph{Current defenses.}
780\raisebox{0.5ex}{+} tests pass; code compiles with \texttt{-Wall
-Werror}; ANSI C99 strict mode with no undefined behavior.

\paragraph{Remaining gaps.}
Fuzzing (AFL\raisebox{0.5ex}{+}), code coverage (gcov/lcov), and memory
sanitizer runs (AddressSanitizer) are not yet formally integrated into CI.

\subsection{Attack Vector 8: Generalization}

\paragraph{Potential challenges.}
Works only for FORTH; results do not generalize; workload too simple;
real-world applicability limited.

\paragraph{Current defenses.}
Workload variety (Fibonacci, Ackermann) and shape-invariance across Zipf
exponents 0.8--1.5 are documented.

\paragraph{Remaining gaps.}
Cross-language validation (Lua, Python) is future work. Real-world benchmark
suites (DaCapo-style) are not tested. Scalability to large programs is
explicitly unvalidated.

\subsection{Attack Vector 9: Conflict of Interest}

\paragraph{Potential challenges.}
Financial interest from patent; results serve IP claims; independent
validation absent; author bias in interpretation.

\paragraph{Current defenses.}
Source code and data are fully public.

\paragraph{Remaining gaps.}
A formal conflict-of-interest statement must appear in all submitted papers.
Template:

\begin{quote}
The author (R.A.\ James) is the named inventor on a provisional patent
application (USPTO, December 2025) related to the adaptive runtime described
herein. To mitigate bias: all data and code are publicly available; statistical
analyses were pre-registered before data collection; independent reproduction
is invited (see \S\ref{sec:replication-invite}). Reviewers may verify all
claims using the provided datasets.
\end{quote}

% PATENT: do not expand or draft claim language beyond what appears here

\subsection{Attack Vector 10: Over-Interpretation}

\paragraph{Potential challenges.}
Too much claimed from limited data; conclusions not supported; speculation
presented as fact; causal claims without causation evidence.

\paragraph{Current defenses.}
Academic wording guidelines (\S\ref{sec:wording-guidelines}) and the ontology
(\S\ref{sec:ontology}) maintain strict language discipline and explicitly
label thermodynamic framing as metaphor.

\paragraph{Remaining gaps.}
Every paper should carry a formal Limitations section distinguishing what the
work demonstrates from what it does not. Speculation must be clearly marked in
running text.

\subsection{Defense Maturity Matrix}
\label{sec:defense-matrix}

\begin{table}[h]
\centering
\caption{Defense readiness by attack vector. Priority: Critical (C),
High (H), Medium (M).}
\label{tab:defense-matrix}
\begin{tabular}{lllll}
\toprule
\textbf{Attack vector} & \textbf{Current} & \textbf{Gap} & \textbf{Priority} \\
\midrule
1. Data integrity & Good & SHA256 + Zenodo & H \\
2. Statistical validity & Partial & Power analysis, R script & C \\
3. Experimental design & Partial & Protocol document & C \\
4. Claims vs.\ evidence & Good & Sensitivity analysis integration & M \\
5. Reproducibility & Partial & Docker + CI/CD & C \\
6. Prior art & Weak & Literature review + table & C \\
7. Implementation & Good & Fuzzing, coverage & M \\
8. Generalization & Weak & Cross-language study & Low \\
9. Conflict of interest & Partial & COI statement in all papers & H \\
10. Over-interpretation & Good & Limitations section everywhere & M \\
\bottomrule
\end{tabular}
\end{table}

\paragraph{Estimated readiness.}
Current state: approximately 65/100. Target for submission: 90/100. After
critical items (statistical R script, protocol document, Docker container,
literature review): approximately 85/100.

Critical bottlenecks are the literature review (time-intensive but essential
for addressing novelty challenges) and independent replication (requires an
external collaborator). Quick wins ($<$ 4 hours each): SHA256 checksums,
conflict-of-interest statement, Docker container, statistical R script.
